\documentclass[11pt,a4paper]{article}
\usepackage[T1]{fontenc}
\usepackage{lmodern}
\usepackage[margin=2.5cm]{geometry}
\usepackage{amsmath,amssymb}
\usepackage{booktabs}
\usepackage{enumitem}
\usepackage{float}
\usepackage{placeins}
\usepackage{microtype}
\usepackage{xcolor}
\usepackage{pgfplots}
\usepackage{pgfplotstable}
\usepgfplotslibrary{groupplots}
\pgfplotsset{compat=1.18}
\usepackage[hidelinks]{hyperref}

\newcommand{\code}[1]{\texttt{#1}}
\setlist{itemsep=2pt,topsep=4pt}
\setlength{\belowcaptionskip}{6pt}

% Chart colours and line styles. Colour: model. Dash: shock.
\definecolor{fcblue}{RGB}{0,84,159}
\definecolor{fcorange}{RGB}{217,95,2}
\pgfplotsset{
  fc/.style={
    xmin=2026, xmax=2075,
    x tick label style={/pgf/number format/1000 sep={}},
    y tick label style={/pgf/number format/fixed},
    scaled y ticks=false,
    tick label style={font=\footnotesize},
    title style={font=\small, yshift=-1ex},
    label style={font=\footnotesize},
    grid=major, grid style={gray!25},
    legend style={font=\footnotesize, draw=none, /tikz/every even column/.append style={column sep=0.8em}},
  },
  narrow/.style={width=5.7cm, height=4.7cm, xtick={2030,2050,2070}},
  wide/.style={width=7.7cm, height=5.4cm, xtick={2030,2040,2050,2060,2070}},
  closureP/.style={fcblue, line width=1.1pt},
  baseP/.style={fcorange, line width=1.1pt},
  closureQ/.style={fcblue, line width=1.1pt, densely dashed},
  baseQ/.style={fcorange, line width=1.1pt, densely dashed},
}

% Four lines of one series. #1: column name without the scenario suffix.
\newcommand{\fourlines}[1]{%
  \addplot[closureP] table[x=year, y=#1FP]{\fcfdata};
  \addplot[baseP] table[x=year, y=#1BP]{\fcfdata};
  \addplot[closureQ] table[x=year, y=#1FQ]{\fcfdata};
  \addplot[baseQ] table[x=year, y=#1BQ]{\fcfdata};
}
\newcommand{\fourlegend}{%
  \legend{{Foreign closure, productivity}, {Base model, productivity},
    {Foreign closure, productivity and fixed cost}, {Base model, productivity and fixed cost}}
}

% Chart data: \fcfdata.
% Chart data for foreign_closure_fixed_cost.tex. A scratch script wrote this file from four solved scenarios.
% Suffix F: foreign closure. B: base model. P: productivity shock. Q: productivity and fixed cost shock.
% Per cent deviation from the baseline, except ecb (per cent).
\pgfplotstableread[row sep=\\]{
year gdpFP gdpBP gdpFQ gdpBQ gvaFP gvaBP gvaFQ gvaBQ conFP conBP conFQ conBQ invFP invBP invFQ invBQ govFP govBP govFQ govBQ expFP expBP expFQ expBQ impFP impBP impFQ impBQ rwFP rwBP rwFQ rwBQ totFP totBP totFQ totBQ pcFP pcBP pcFQ pcBQ ecbFP ecbBP ecbFQ ecbBQ capFP capBP capFQ capBQ matFP matBP matFQ matBQ bunFP bunBP bunFQ bunBQ outFP outBP outFQ outBQ \\
2026 0.58606 0.61870 0.34193 0.39976 0.65994 0.65806 0.43215 0.44986 0.06993 0.31236 -0.27283 0.11292 1.90923 2.43071 1.78830 2.72465 0.76967 0.79550 0.64264 0.67529 0.58703 0.09696 0.42137 -0.13076 0.61920 0.47666 0.45627 0.48365 0.41238 0.35805 0.11418 0.19482 0.00122 -0.01101 -0.01228 0.04301 -0.00000 0.10512 -0.00000 0.38165 4.00000 4.00000 4.00000 4.00000 0.13947 0.15901 0.11458 0.14423 0.60843 0.56814 0.55138 0.59949 0.55389 0.53577 0.52831 0.55127 0.63024 0.60631 0.47997 0.50404 \\
2027 0.65894 0.74141 0.40424 0.51089 0.75778 0.77304 0.52235 0.56386 -0.00556 0.48144 -0.36942 0.24593 2.46578 2.92943 2.46442 3.81275 0.83433 0.88074 0.62717 0.70331 0.68577 -0.03518 0.49306 -0.41744 0.73139 0.49363 0.54146 0.49930 0.51642 0.56814 0.32772 0.53719 0.00770 0.01214 -0.00114 0.09246 -0.00000 0.08442 -0.00000 0.22853 3.61067 4.00000 4.17061 4.00000 0.28551 0.30653 0.24632 0.32863 0.72403 0.66021 0.68797 0.69492 0.64442 0.62039 0.62162 0.63309 0.73455 0.70389 0.58755 0.59887 \\
2028 0.75478 0.82336 0.48457 0.60153 0.85271 0.84191 0.60308 0.64427 0.08864 0.60173 -0.27886 0.37010 2.61006 2.26312 2.56805 3.02647 0.89031 0.92302 0.67494 0.75150 0.75572 0.18863 0.57563 -0.14164 0.79441 0.49159 0.62323 0.50204 0.67236 0.61265 0.47092 0.59145 0.00197 -0.03539 -0.00202 0.03438 0.00000 0.00501 0.00000 0.10707 4.31930 4.00000 4.47552 4.00000 0.41521 0.41445 0.37210 0.47863 0.82117 0.67619 0.76100 0.69213 0.72792 0.66669 0.69017 0.68039 0.82851 0.75356 0.66463 0.64876 \\
2029 0.84162 0.88511 0.56703 0.68445 0.92880 0.89413 0.67478 0.71661 0.24141 0.68334 -0.12282 0.46741 2.53812 1.77593 2.47438 2.38574 0.94276 0.95422 0.72535 0.79356 0.82101 0.35944 0.64643 0.09685 0.85180 0.48760 0.68814 0.49851 0.78904 0.64625 0.58275 0.64041 -0.00106 -0.07171 -0.00321 -0.01642 -0.00000 -0.04975 -0.00000 0.02603 4.24505 4.00000 4.36693 4.00000 0.52809 0.49001 0.48396 0.58754 0.88300 0.69498 0.81468 0.71698 0.79075 0.70386 0.74824 0.73148 0.89806 0.79381 0.72875 0.70408 \\
2030 0.91768 0.92830 0.64194 0.74622 0.99124 0.93113 0.73602 0.77041 0.40514 0.73520 0.04301 0.53654 2.38609 1.43745 2.31477 1.92088 0.98945 0.97582 0.77170 0.82459 0.88144 0.48240 0.70899 0.27296 0.90536 0.48361 0.74440 0.49431 0.87942 0.67007 0.67182 0.67687 -0.00208 -0.09780 -0.00365 -0.05392 0.00000 -0.08834 0.00000 -0.03003 4.17229 4.00000 4.28603 4.00000 0.62458 0.54227 0.58037 0.66420 0.92626 0.70971 0.85541 0.73968 0.84054 0.73062 0.79664 0.77074 0.95260 0.82289 0.78170 0.74682 \\
2031 0.98297 0.95800 0.70711 0.79007 1.04298 0.95682 0.78765 0.80849 0.55841 0.76834 0.19737 0.58567 2.21680 1.20290 2.14427 1.59114 1.02982 0.99058 0.81223 0.84660 0.93444 0.56935 0.76302 0.39750 0.95249 0.47998 0.79268 0.49048 0.95162 0.68645 0.74413 0.70249 -0.00228 -0.11619 -0.00368 -0.08038 0.00000 -0.11562 0.00000 -0.06913 4.11354 4.00000 4.22464 4.00000 0.70615 0.57818 0.66225 0.71751 0.95870 0.72013 0.88730 0.75616 0.88101 0.74924 0.83688 0.79862 0.99664 0.84316 0.82547 0.77720 \\
2032 1.03833 0.97822 0.76269 0.82067 1.08597 0.97448 0.83095 0.83497 0.69388 0.78937 0.33337 0.62043 2.05565 1.04030 1.98482 1.35915 1.06412 1.00058 0.84686 0.86200 0.97972 0.63043 0.80903 0.48423 0.99280 0.47678 0.83371 0.48720 1.01027 0.69750 0.80353 0.72012 -0.00216 -0.12905 -0.00351 -0.09874 -0.00000 -0.13497 -0.00000 -0.09651 4.06592 4.00000 4.17594 4.00000 0.77463 0.60270 0.73127 0.75434 0.98409 0.72719 0.91289 0.76735 0.91428 0.76199 0.87038 0.81790 1.03271 0.85707 0.86179 0.79820 \\
2033 1.08490 0.99187 0.80964 0.84186 1.12170 0.98652 0.86720 0.85325 0.81032 0.80244 0.45011 0.64493 1.91150 0.92749 1.84377 1.19649 1.09299 1.00728 0.87614 0.87270 1.01788 0.67323 0.84788 0.54421 1.02679 0.47399 0.86835 0.48442 1.05829 0.70481 0.85268 0.73211 -0.00191 -0.13801 -0.00326 -0.11139 0.00000 -0.14873 0.00000 -0.11571 4.02703 4.00000 4.13637 4.00000 0.83184 0.61932 0.78922 0.77965 1.00441 0.73183 0.93377 0.77472 0.94176 0.77062 0.89829 0.83106 1.06242 0.86649 0.89201 0.81255 \\
2034 1.12386 1.00098 0.84912 0.85644 1.15136 0.99465 0.89750 0.86579 0.90898 0.81024 0.54903 0.66213 1.78614 0.84920 1.72218 1.08266 1.11714 1.01172 0.90075 0.88010 1.04978 0.70320 0.88053 0.58548 1.05522 0.47157 0.89749 0.48210 1.09771 0.70952 0.89349 0.74019 -0.00162 -0.14423 -0.00300 -0.12004 0.00000 -0.15855 0.00000 -0.12918 3.99512 4.00000 4.10382 4.00000 0.87945 0.63049 0.83776 0.79695 1.02086 0.73480 0.95097 0.77947 0.96447 0.77640 0.92157 0.83997 1.08694 0.87281 0.91718 0.82226 \\
2035 1.15635 1.00696 0.88221 0.86642 1.17591 1.00009 0.92283 0.87435 0.99190 0.81451 0.63228 0.67416 1.67857 0.79487 1.61871 1.00308 1.13724 1.01461 0.92136 0.88520 1.07630 0.72420 0.90789 0.61373 1.07886 0.46948 0.92193 0.48016 1.13004 0.71245 0.92742 0.74556 -0.00135 -0.14855 -0.00274 -0.12593 0.00000 -0.16561 0.00000 -0.13865 3.96888 4.00000 4.07689 4.00000 0.91893 0.63789 0.87835 0.80869 1.03422 0.73664 0.96520 0.78246 0.98323 0.78020 0.94100 0.84595 1.10716 0.87698 0.93817 0.82877 \\
2036 1.18333 1.01079 0.90991 0.87320 1.19619 1.00366 0.94396 0.88013 1.06126 0.81645 0.70209 0.68254 1.58684 0.75716 1.53119 0.94752 1.15392 1.01642 0.93860 0.88869 1.09824 0.73894 0.93077 0.63296 1.09843 0.46767 0.94238 0.47856 1.15648 0.71416 0.95563 0.74908 -0.00108 -0.15155 -0.00249 -0.12989 0.00000 -0.17071 0.00000 -0.14530 3.94732 4.00000 4.05455 4.00000 0.95155 0.64270 0.91224 0.81658 1.04506 0.73771 0.97700 0.78427 0.99867 0.78266 0.95720 0.84992 1.12379 0.87968 0.95567 0.83309 \\
2037 1.20568 1.01315 0.93306 0.87775 1.21287 1.00596 0.96159 0.88401 1.11908 0.81683 0.76051 0.68837 1.50884 0.73099 1.45740 0.90877 1.16771 1.01750 0.95300 0.89107 1.11632 0.74933 0.94987 0.64594 1.11455 0.46609 0.95946 0.47722 1.17800 0.71505 0.97907 0.75132 -0.00084 -0.15362 -0.00227 -0.13253 0.00000 -0.17444 0.00000 -0.14999 3.92963 4.00000 4.03601 4.00000 0.97838 0.64574 0.94050 0.82182 1.05380 0.73827 0.98679 0.78531 1.01133 0.78419 0.97070 0.85252 1.13742 0.88138 0.97025 0.83591 \\
2038 1.22413 1.01450 0.95239 0.88077 1.22654 1.00738 0.97627 0.88657 1.16712 0.81621 0.80933 0.69239 1.44265 0.71278 1.39526 0.88176 1.17908 1.01809 0.96500 0.89266 1.13115 0.75671 0.96579 0.65462 1.12776 0.46471 0.97372 0.47611 1.19541 0.71538 0.99851 0.75270 -0.00062 -0.15506 -0.00207 -0.13427 0.00000 -0.17719 0.00000 -0.15329 3.91519 4.00000 4.02062 4.00000 1.00035 0.64756 0.96403 0.82524 1.06082 0.73851 0.99490 0.78585 1.02167 0.78510 0.98194 0.85418 1.14853 0.88241 0.98237 0.83771 \\
2039 1.23929 1.01518 0.96850 0.88273 1.23769 1.00821 0.98848 0.88823 1.20689 0.81495 0.85009 0.69514 1.38665 0.70000 1.34301 0.86293 1.18841 1.01834 0.97500 0.89372 1.14325 0.76202 0.97905 0.66037 1.13856 0.46350 0.98559 0.47518 1.20942 0.71534 1.01462 0.75348 -0.00042 -0.15607 -0.00189 -0.13539 0.00000 -0.17924 0.00000 -0.15562 3.90347 4.00000 4.00786 4.00000 1.01825 0.64855 0.98359 0.82741 1.06640 0.73852 1.00161 0.78606 1.03006 0.78559 0.99128 0.85521 1.15755 0.88297 0.99245 0.83882 \\
2040 1.25169 1.01542 0.98192 0.88398 1.24675 1.00864 0.99863 0.88927 1.23963 0.81326 0.88399 0.69697 1.33933 0.69124 1.29878 0.84976 1.19604 1.01837 0.98332 0.89440 1.15308 0.76591 0.99005 0.66411 1.14734 0.46243 0.99547 0.47441 1.22066 0.71507 1.02795 0.75387 -0.00024 -0.15679 -0.00173 -0.13610 0.00000 -0.18081 0.00000 -0.15727 3.89406 4.00000 3.99730 4.00000 1.03280 0.64900 0.99985 0.82873 1.07084 0.73843 1.00716 0.78608 1.03686 0.78580 0.99904 0.85581 1.16485 0.88324 1.00081 0.83948 \\
2041 1.26179 1.01537 0.99308 0.88474 1.25406 1.00882 1.00705 0.88990 1.26650 0.81133 0.91219 0.69815 1.29954 0.68527 1.26152 0.84060 1.20225 1.01825 0.99025 0.89483 1.16103 0.76881 0.99917 0.66649 1.15443 0.46149 1.00368 0.47377 1.22961 0.71465 1.03897 0.75400 -0.00009 -0.15730 -0.00159 -0.13652 -0.00000 -0.18203 -0.00000 -0.15844 3.88649 4.00000 3.98858 4.00000 1.04454 0.64910 1.01333 0.82948 1.07435 0.73827 1.01174 0.78598 1.04234 0.78585 1.00547 0.85614 1.17071 0.88332 1.00775 0.83982 \\
2042 1.26997 1.01513 1.00235 0.88516 1.25993 1.00881 1.01402 0.89024 1.28850 0.80929 0.93564 0.69889 1.26617 0.68122 1.23021 0.83426 1.20727 1.01804 0.99599 0.89509 1.16741 0.77102 1.00672 0.66795 1.16013 0.46065 1.01049 0.47323 1.23668 0.71415 1.04805 0.75396 0.00005 -0.15768 -0.00147 -0.13675 0.00000 -0.18300 0.00000 -0.15927 3.88043 4.00000 3.98139 4.00000 1.05397 0.64897 1.02449 0.82985 1.07708 0.73807 1.01553 0.78582 1.04672 0.78578 1.01081 0.85628 1.17539 0.88328 1.01348 0.83996 \\
2043 1.27656 1.01477 1.01003 0.88536 1.26460 1.00870 1.01979 0.89040 1.30642 0.80720 0.95513 0.69932 1.23824 0.67851 1.20394 0.82990 1.21132 1.01777 1.00076 0.89522 1.17249 0.77275 1.01297 0.66879 1.16468 0.45990 1.01613 0.47278 1.24219 0.71361 1.05553 0.75381 0.00017 -0.15795 -0.00135 -0.13686 0.00000 -0.18378 0.00000 -0.15986 3.87562 4.00000 3.97546 4.00000 1.06148 0.64869 1.03372 0.82996 1.07918 0.73787 1.01864 0.78562 1.05020 0.78565 1.01521 0.85630 1.17909 0.88315 1.01823 0.83998 \\
2044 1.28182 1.01434 1.01640 0.88542 1.26829 1.00851 1.02455 0.89043 1.32098 0.80512 0.97131 0.69954 1.21489 0.67672 1.18192 0.82692 1.21455 1.01747 1.00470 0.89528 1.17651 0.77414 1.01813 0.66922 1.16826 0.45923 1.02079 0.47240 1.24642 0.71305 1.06167 0.75360 0.00028 -0.15816 -0.00126 -0.13689 0.00000 -0.18443 0.00000 -0.16029 3.87184 4.00000 3.97059 4.00000 1.06741 0.64833 1.04133 0.82991 1.08076 0.73766 1.02119 0.78541 1.05292 0.78549 1.01885 0.85625 1.18198 0.88299 1.02214 0.83992 \\
2045 1.28600 1.01388 1.02165 0.88538 1.27117 1.00829 1.02847 0.89039 1.33273 0.80309 0.98473 0.69961 1.19538 0.67557 1.16347 0.82489 1.21710 1.01716 1.00796 0.89528 1.17966 0.77529 1.02239 0.66940 1.17106 0.45862 1.02463 0.47209 1.24961 0.71249 1.06669 0.75335 0.00037 -0.15833 -0.00117 -0.13686 -0.00000 -0.18498 -0.00000 -0.16061 3.86890 4.00000 3.96658 4.00000 1.07204 0.64792 1.04759 0.82977 1.08192 0.73746 1.02328 0.78521 1.05504 0.78530 1.02184 0.85617 1.18422 0.88280 1.02535 0.83982 \\
2046 1.28928 1.01340 1.02599 0.88527 1.27337 1.00805 1.03169 0.89030 1.34217 0.80112 0.99584 0.69958 1.17911 0.67486 1.14802 0.82354 1.21911 1.01684 1.01065 0.89526 1.18208 0.77628 1.02588 0.66941 1.17321 0.45807 1.02779 0.47182 1.25195 0.71194 1.07079 0.75308 0.00045 -0.15846 -0.00109 -0.13680 0.00000 -0.18546 0.00000 -0.16084 3.86665 4.00000 3.96331 4.00000 1.07559 0.64750 1.05273 0.82956 1.08274 0.73728 1.02498 0.78503 1.05664 0.78511 1.02430 0.85605 1.18591 0.88261 1.02799 0.83969 \\
2047 1.29183 1.01292 1.02955 0.88514 1.27504 1.00780 1.03433 0.89019 1.34970 0.79923 1.00503 0.69949 1.16554 0.67444 1.13509 0.82264 1.22066 1.01652 1.01286 0.89521 1.18391 0.77715 1.02875 0.66933 1.17483 0.45757 1.03039 0.47159 1.25360 0.71141 1.07411 0.75281 0.00051 -0.15858 -0.00103 -0.13673 -0.00000 -0.18588 -0.00000 -0.16102 3.86496 4.00000 3.96064 4.00000 1.07827 0.64707 1.05694 0.82933 1.08328 0.73712 1.02636 0.78485 1.05784 0.78492 1.02631 0.85593 1.18716 0.88241 1.03014 0.83955 \\
2048 1.29376 1.01246 1.03248 0.88498 1.27626 1.00755 1.03648 0.89006 1.35566 0.79742 1.01261 0.69935 1.15424 0.67422 1.12428 0.82206 1.22183 1.01622 1.01468 0.89515 1.18526 0.77793 1.03110 0.66919 1.17602 0.45711 1.03251 0.47140 1.25469 0.71090 1.07680 0.75254 0.00057 -0.15867 -0.00097 -0.13665 0.00000 -0.18626 0.00000 -0.16116 3.86372 4.00000 3.95846 4.00000 1.08024 0.64665 1.06037 0.82908 1.08359 0.73696 1.02748 0.78470 1.05870 0.78474 1.02795 0.85581 1.18804 0.88223 1.03190 0.83941 \\
2049 1.29521 1.01200 1.03487 0.88482 1.27711 1.00731 1.03823 0.88992 1.36031 0.79570 1.01887 0.69918 1.14483 0.67414 1.11525 0.82170 1.22270 1.01592 1.01616 0.89508 1.18622 0.77865 1.03301 0.66902 1.17685 0.45669 1.03423 0.47124 1.25533 0.71042 1.07896 0.75228 0.00062 -0.15876 -0.00092 -0.13656 0.00000 -0.18660 0.00000 -0.16127 3.86285 4.00000 3.95670 4.00000 1.08163 0.64625 1.06316 0.82883 1.08373 0.73683 1.02838 0.78456 1.05929 0.78456 1.02928 0.85570 1.18864 0.88205 1.03333 0.83927 \\
2050 1.29625 1.01157 1.03682 0.88465 1.27768 1.00707 1.03965 0.88979 1.36391 0.79407 1.02401 0.69899 1.13702 0.67416 1.10770 0.82149 1.22333 1.01564 1.01737 0.89501 1.18686 0.77932 1.03456 0.66885 1.17740 0.45631 1.03563 0.47109 1.25562 0.70996 1.08068 0.75203 0.00066 -0.15884 -0.00088 -0.13648 0.00000 -0.18692 0.00000 -0.16136 3.86228 4.00000 3.95528 4.00000 1.08256 0.64586 1.06541 0.82858 1.08372 0.73671 1.02909 0.78444 1.05966 0.78440 1.03036 0.85559 1.18901 0.88188 1.03447 0.83915 \\
2051 1.29696 1.01116 1.03840 0.88449 1.27801 1.00685 1.04079 0.88966 1.36664 0.79252 1.02823 0.69879 1.13053 0.67423 1.10141 0.82138 1.22375 1.01537 1.01835 0.89494 1.18725 0.77995 1.03581 0.66868 1.17772 0.45595 1.03676 0.47097 1.25562 0.70953 1.08204 0.75179 0.00070 -0.15892 -0.00084 -0.13640 0.00000 -0.18722 0.00000 -0.16144 3.86195 4.00000 3.95414 4.00000 1.08310 0.64550 1.06722 0.82835 1.08361 0.73660 1.02966 0.78434 1.05986 0.78425 1.03122 0.85548 1.18918 0.88172 1.03540 0.83903 \\
2052 1.29741 1.01077 1.03968 0.88433 1.27816 1.00665 1.04170 0.88954 1.36865 0.79105 1.03168 0.69859 1.12515 0.67434 1.09615 0.82133 1.22401 1.01511 1.01915 0.89487 1.18743 0.78054 1.03681 0.66853 1.17786 0.45562 1.03766 0.47086 1.25540 0.70912 1.08311 0.75157 0.00073 -0.15899 -0.00081 -0.13632 0.00000 -0.18749 0.00000 -0.16151 3.86180 4.00000 3.95323 4.00000 1.08335 0.64516 1.06866 0.82813 1.08341 0.73650 1.03010 0.78424 1.05991 0.78411 1.03191 0.85539 1.18921 0.88158 1.03613 0.83892 \\
2053 1.29766 1.01041 1.04071 0.88418 1.27816 1.00645 1.04242 0.88942 1.37010 0.78966 1.03449 0.69839 1.12069 0.67447 1.09178 0.82133 1.22415 1.01487 1.01979 0.89480 1.18746 0.78111 1.03760 0.66838 1.17785 0.45532 1.03838 0.47076 1.25501 0.70873 1.08393 0.75136 0.00076 -0.15905 -0.00078 -0.13625 -0.00000 -0.18775 -0.00000 -0.16157 3.86181 4.00000 3.95251 4.00000 1.08336 0.64484 1.06981 0.82792 1.08315 0.73642 1.03044 0.78416 1.05986 0.78398 1.03246 0.85530 1.18912 0.88144 1.03671 0.83882 \\
2054 1.29773 1.01006 1.04153 0.88404 1.27804 1.00627 1.04299 0.88932 1.37108 0.78835 1.03678 0.69819 1.11700 0.67462 1.08813 0.82136 1.22418 1.01464 1.02029 0.89473 1.18736 0.78165 1.03823 0.66826 1.17773 0.45503 1.03894 0.47068 1.25449 0.70837 1.08456 0.75116 0.00078 -0.15912 -0.00076 -0.13619 0.00000 -0.18800 0.00000 -0.16162 3.86193 4.00000 3.95195 4.00000 1.08319 0.64454 1.07070 0.82772 1.08284 0.73634 1.03069 0.78410 1.05972 0.78386 1.03288 0.85522 1.18894 0.88132 1.03716 0.83873 \\
2055 1.29768 1.00974 1.04217 0.88391 1.27784 1.00610 1.04343 0.88922 1.37169 0.78711 1.03863 0.69799 1.11396 0.67477 1.08510 0.82141 1.22413 1.01442 1.02070 0.89467 1.18716 0.78217 1.03872 0.66814 1.17753 0.45477 1.03938 0.47060 1.25388 0.70803 1.08502 0.75097 0.00080 -0.15918 -0.00074 -0.13613 0.00000 -0.18823 0.00000 -0.16168 3.86213 4.00000 3.95152 4.00000 1.08287 0.64426 1.07140 0.82754 1.08249 0.73627 1.03088 0.78404 1.05951 0.78375 1.03321 0.85515 1.18870 0.88121 1.03751 0.83865 \\
2056 1.29752 1.00943 1.04268 0.88378 1.27757 1.00594 1.04377 0.88912 1.37200 0.78594 1.04012 0.69779 1.11144 0.67493 1.08258 0.82147 1.22402 1.01422 1.02101 0.89461 1.18688 0.78266 1.03910 0.66805 1.17725 0.45452 1.03971 0.47054 1.25320 0.70771 1.08535 0.75080 0.00082 -0.15924 -0.00072 -0.13607 0.00000 -0.18844 0.00000 -0.16173 3.86240 4.00000 3.95119 4.00000 1.08245 0.64400 1.07192 0.82737 1.08212 0.73621 1.03101 0.78399 1.05926 0.78365 1.03346 0.85509 1.18840 0.88110 1.03777 0.83857 \\
2057 1.29728 1.00914 1.04307 0.88367 1.27725 1.00579 1.04402 0.88904 1.37207 0.78483 1.04131 0.69760 1.10937 0.67509 1.08048 0.82154 1.22386 1.01403 1.02126 0.89456 1.18655 0.78313 1.03939 0.66796 1.17693 0.45429 1.03997 0.47048 1.25247 0.70742 1.08558 0.75064 0.00083 -0.15930 -0.00070 -0.13602 -0.00000 -0.18865 -0.00000 -0.16178 3.86273 4.00000 3.95094 4.00000 1.08194 0.64376 1.07231 0.82722 1.08173 0.73616 1.03110 0.78394 1.05897 0.78355 1.03364 0.85503 1.18807 0.88101 1.03796 0.83851 \\
2058 1.29699 1.00888 1.04336 0.88356 1.27689 1.00565 1.04420 0.88896 1.37196 0.78379 1.04226 0.69742 1.10766 0.67526 1.07874 0.82161 1.22367 1.01385 1.02144 0.89450 1.18618 0.78358 1.03959 0.66789 1.17657 0.45408 1.04015 0.47042 1.25172 0.70714 1.08572 0.75049 0.00084 -0.15935 -0.00069 -0.13598 0.00000 -0.18885 0.00000 -0.16182 3.86308 4.00000 3.95076 4.00000 1.08138 0.64353 1.07260 0.82707 1.08134 0.73612 1.03116 0.78391 1.05866 0.78347 1.03377 0.85498 1.18771 0.88092 1.03809 0.83845 \\
2059 1.29665 1.00863 1.04358 0.88345 1.27651 1.00553 1.04433 0.88888 1.37171 0.78282 1.04300 0.69724 1.10626 0.67542 1.07730 0.82169 1.22345 1.01368 1.02158 0.89445 1.18579 0.78400 1.03974 0.66783 1.17619 0.45388 1.04027 0.47037 1.25095 0.70688 1.08579 0.75035 0.00085 -0.15940 -0.00068 -0.13594 -0.00000 -0.18903 -0.00000 -0.16187 3.86346 4.00000 3.95064 4.00000 1.08077 0.64333 1.07279 0.82694 1.08095 0.73608 1.03118 0.78387 1.05832 0.78339 1.03386 0.85493 1.18733 0.88084 1.03818 0.83839 \\
2060 1.29628 1.00839 1.04373 0.88336 1.27611 1.00541 1.04441 0.88882 1.37135 0.78190 1.04358 0.69707 1.10512 0.67559 1.07611 0.82176 1.22321 1.01353 1.02168 0.89441 1.18537 0.78440 1.03984 0.66778 1.17580 0.45369 1.04036 0.47032 1.25018 0.70664 1.08581 0.75022 0.00086 -0.15945 -0.00067 -0.13590 0.00000 -0.18920 0.00000 -0.16191 3.86385 4.00000 3.95056 4.00000 1.08015 0.64313 1.07291 0.82682 1.08056 0.73604 1.03118 0.78385 1.05798 0.78332 1.03392 0.85489 1.18695 0.88077 1.03823 0.83834 \\
2061 1.29590 1.00817 1.04384 0.88327 1.27571 1.00531 1.04446 0.88875 1.37091 0.78103 1.04403 0.69690 1.10418 0.67576 1.07513 0.82184 1.22297 1.01338 1.02175 0.89436 1.18495 0.78478 1.03990 0.66774 1.17540 0.45352 1.04040 0.47028 1.24941 0.70642 1.08578 0.75010 0.00086 -0.15950 -0.00066 -0.13586 0.00000 -0.18937 0.00000 -0.16195 3.86425 4.00000 3.95052 4.00000 1.07951 0.64296 1.07297 0.82670 1.08017 0.73601 1.03116 0.78382 1.05764 0.78325 1.03394 0.85485 1.18656 0.88071 1.03826 0.83830 \\
2062 1.29550 1.00797 1.04390 0.88318 1.27531 1.00521 1.04447 0.88870 1.37041 0.78023 1.04436 0.69674 1.10343 0.67594 1.07432 0.82192 1.22271 1.01325 1.02179 0.89432 1.18453 0.78514 1.03992 0.66770 1.17500 0.45336 1.04042 0.47024 1.24866 0.70622 1.08573 0.74999 0.00087 -0.15955 -0.00066 -0.13583 0.00000 -0.18952 0.00000 -0.16199 3.86465 4.00000 3.95051 4.00000 1.07887 0.64280 1.07299 0.82660 1.07980 0.73599 1.03113 0.78380 1.05730 0.78319 1.03395 0.85482 1.18618 0.88065 1.03826 0.83826 \\
2063 1.29510 1.00779 1.04394 0.88311 1.27491 1.00512 1.04447 0.88864 1.36987 0.77947 1.04460 0.69659 1.10282 0.67612 1.07366 0.82200 1.22246 1.01312 1.02181 0.89429 1.18411 0.78548 1.03993 0.66767 1.17460 0.45321 1.04041 0.47021 1.24793 0.70604 1.08565 0.74988 0.00088 -0.15959 -0.00065 -0.13581 -0.00000 -0.18967 -0.00000 -0.16203 3.86505 4.00000 3.95051 4.00000 1.07823 0.64265 1.07297 0.82651 1.07944 0.73597 1.03109 0.78379 1.05696 0.78314 1.03394 0.85479 1.18580 0.88060 1.03824 0.83822 \\
2064 1.29471 1.00761 1.04394 0.88304 1.27452 1.00504 1.04444 0.88860 1.36930 0.77877 1.04477 0.69645 1.10235 0.67631 1.07312 0.82209 1.22221 1.01301 1.02182 0.89425 1.18370 0.78579 1.03991 0.66765 1.17421 0.45308 1.04039 0.47018 1.24722 0.70587 1.08555 0.74979 0.00088 -0.15963 -0.00065 -0.13578 -0.00000 -0.18980 -0.00000 -0.16206 3.86544 4.00000 3.95054 4.00000 1.07761 0.64252 1.07292 0.82642 1.07909 0.73596 1.03105 0.78378 1.05663 0.78309 1.03391 0.85476 1.18543 0.88056 1.03821 0.83819 \\
2065 1.29432 1.00746 1.04393 0.88297 1.27414 1.00497 1.04441 0.88855 1.36871 0.77812 1.04488 0.69632 1.10201 0.67650 1.07269 0.82218 1.22196 1.01290 1.02181 0.89422 1.18330 0.78608 1.03988 0.66763 1.17383 0.45296 1.04036 0.47015 1.24655 0.70572 1.08544 0.74970 0.00088 -0.15966 -0.00064 -0.13576 0.00000 -0.18993 0.00000 -0.16210 3.86581 4.00000 3.95057 4.00000 1.07700 0.64239 1.07286 0.82634 1.07875 0.73595 1.03100 0.78377 1.05631 0.78305 1.03388 0.85474 1.18508 0.88052 1.03817 0.83816 \\
2066 1.29394 1.00732 1.04390 0.88291 1.27378 1.00491 1.04436 0.88851 1.36811 0.77751 1.04494 0.69619 1.10178 0.67670 1.07234 0.82227 1.22172 1.01281 1.02179 0.89419 1.18292 0.78635 1.03984 0.66761 1.17348 0.45285 1.04031 0.47012 1.24591 0.70558 1.08532 0.74963 0.00089 -0.15970 -0.00064 -0.13574 -0.00000 -0.19005 -0.00000 -0.16213 3.86618 4.00000 3.95062 4.00000 1.07642 0.64229 1.07278 0.82627 1.07844 0.73594 1.03094 0.78376 1.05601 0.78301 1.03384 0.85472 1.18474 0.88049 1.03812 0.83814 \\
2067 1.29357 1.00719 1.04386 0.88285 1.27344 1.00485 1.04431 0.88848 1.36750 0.77696 1.04496 0.69607 1.10168 0.67691 1.07209 0.82237 1.22149 1.01272 1.02177 0.89416 1.18255 0.78659 1.03979 0.66759 1.17314 0.45275 1.04026 0.47010 1.24531 0.70547 1.08520 0.74956 0.00089 -0.15973 -0.00063 -0.13572 0.00000 -0.19015 0.00000 -0.16215 3.86654 4.00000 3.95067 4.00000 1.07587 0.64219 1.07269 0.82621 1.07815 0.73594 1.03089 0.78375 1.05572 0.78298 1.03379 0.85470 1.18442 0.88046 1.03807 0.83812 \\
2068 1.29322 1.00707 1.04381 0.88280 1.27312 1.00481 1.04426 0.88845 1.36689 0.77646 1.04494 0.69596 1.10172 0.67713 1.07191 0.82248 1.22127 1.01265 1.02174 0.89413 1.18221 0.78680 1.03974 0.66758 1.17282 0.45266 1.04021 0.47007 1.24476 0.70537 1.08508 0.74950 0.00089 -0.15975 -0.00063 -0.13570 -0.00000 -0.19025 -0.00000 -0.16218 3.86688 4.00000 3.95073 4.00000 1.07535 0.64211 1.07260 0.82615 1.07788 0.73594 1.03083 0.78375 1.05545 0.78295 1.03375 0.85468 1.18413 0.88044 1.03802 0.83810 \\
2069 1.29289 1.00697 1.04376 0.88276 1.27283 1.00477 1.04420 0.88842 1.36626 0.77601 1.04490 0.69586 1.10192 0.67735 1.07180 0.82259 1.22106 1.01258 1.02171 0.89411 1.18190 0.78700 1.03968 0.66756 1.17254 0.45259 1.04016 0.47006 1.24427 0.70529 1.08497 0.74945 0.00090 -0.15978 -0.00063 -0.13568 0.00000 -0.19034 0.00000 -0.16220 3.86720 4.00000 3.95079 4.00000 1.07487 0.64204 1.07251 0.82610 1.07765 0.73594 1.03078 0.78375 1.05520 0.78293 1.03370 0.85467 1.18386 0.88042 1.03796 0.83808 \\
2070 1.29258 1.00689 1.04371 0.88272 1.27256 1.00474 1.04415 0.88840 1.36562 0.77561 1.04483 0.69577 1.10233 0.67759 1.07179 0.82271 1.22087 1.01252 1.02168 0.89409 1.18163 0.78716 1.03963 0.66755 1.17229 0.45253 1.04011 0.47004 1.24385 0.70522 1.08487 0.74941 0.00090 -0.15980 -0.00063 -0.13567 -0.00000 -0.19042 -0.00000 -0.16222 3.86750 4.00000 3.95085 4.00000 1.07445 0.64198 1.07242 0.82606 1.07745 0.73595 1.03074 0.78375 1.05498 0.78291 1.03366 0.85466 1.18363 0.88042 1.03792 0.83807 \\
2071 1.29230 1.00682 1.04366 0.88268 1.27234 1.00472 1.04410 0.88838 1.36496 0.77526 1.04473 0.69569 1.10300 0.67783 1.07186 0.82284 1.22069 1.01248 1.02165 0.89407 1.18140 0.78730 1.03959 0.66753 1.17209 0.45248 1.04007 0.47003 1.24352 0.70518 1.08479 0.74937 0.00090 -0.15982 -0.00062 -0.13566 -0.00000 -0.19049 -0.00000 -0.16224 3.86779 4.00000 3.95091 4.00000 1.07409 0.64194 1.07235 0.82603 1.07730 0.73596 1.03071 0.78375 1.05480 0.78290 1.03362 0.85465 1.18343 0.88041 1.03788 0.83806 \\
2072 1.29205 1.00676 1.04360 0.88265 1.27217 1.00471 1.04407 0.88836 1.36425 0.77497 1.04461 0.69562 1.10400 0.67807 1.07205 0.82299 1.22053 1.01244 1.02162 0.89406 1.18122 0.78742 1.03956 0.66751 1.17195 0.45244 1.04005 0.47002 1.24329 0.70514 1.08473 0.74935 0.00090 -0.15983 -0.00062 -0.13565 0.00000 -0.19055 0.00000 -0.16225 3.86805 4.00000 3.95097 4.00000 1.07382 0.64191 1.07229 0.82600 1.07720 0.73597 1.03069 0.78375 1.05466 0.78289 1.03359 0.85465 1.18329 0.88041 1.03785 0.83805 \\
2073 1.29184 1.00672 1.04356 0.88263 1.27206 1.00470 1.04404 0.88835 1.36349 0.77473 1.04445 0.69556 1.10546 0.67831 1.07238 0.82315 1.22040 1.01242 1.02159 0.89405 1.18112 0.78750 1.03954 0.66748 1.17188 0.45241 1.04003 0.47002 1.24320 0.70513 1.08471 0.74933 0.00090 -0.15984 -0.00062 -0.13563 0.00000 -0.19059 0.00000 -0.16226 3.86829 4.00000 3.95102 4.00000 1.07365 0.64188 1.07225 0.82598 1.07718 0.73599 1.03068 0.78376 1.05458 0.78289 1.03358 0.85464 1.18322 0.88042 1.03783 0.83805 \\
2074 1.29168 1.00669 1.04352 0.88262 1.27202 1.00470 1.04404 0.88835 1.36264 0.77455 1.04426 0.69552 1.10751 0.67854 1.07288 0.82334 1.22029 1.01240 1.02156 0.89404 1.18110 0.78756 1.03954 0.66745 1.17192 0.45239 1.04005 0.47003 1.24328 0.70512 1.08473 0.74933 0.00089 -0.15985 -0.00062 -0.13562 -0.00000 -0.19063 -0.00000 -0.16227 3.86849 4.00000 3.95107 4.00000 1.07361 0.64187 1.07225 0.82597 1.07724 0.73601 1.03070 0.78377 1.05456 0.78289 1.03358 0.85464 1.18322 0.88043 1.03783 0.83805 \\
2075 1.29166 1.00668 1.04352 0.88261 1.27207 1.00471 1.04405 0.88835 1.36214 0.77444 1.04415 0.69549 1.10909 0.67874 1.07326 0.82355 1.22021 1.01239 1.02155 0.89404 1.18109 0.78758 1.03954 0.66740 1.17193 0.45238 1.04005 0.47004 1.24357 0.70513 1.08480 0.74935 0.00088 -0.15985 -0.00063 -0.13561 0.00053 -0.19066 0.00014 -0.16227 3.87059 4.00000 3.95163 4.00000 1.07361 0.64187 1.07225 0.82598 1.07744 0.73603 1.03076 0.78379 1.05464 0.78289 1.03360 0.85465 1.18332 0.88044 1.03786 0.83806 \\
}\fcfdata


\title{The foreign closure:\\ why GDP rises more than one per cent, and a shock to fixed costs}
\author{}
\date{27 November 2026}

\begin{document}
\maketitle

\section{Summary}

The note \code{foreign\_closure.tex} shows a shock of one per cent to labour productivity.
In 2075, GDP is 1.29 per cent above the baseline with the foreign closure and 1.01 per cent above in the base model.
This note gives the cause of both numbers and tests it with a second shock.

\begin{enumerate}
\item \textbf{Employment does not change.} Labour is 27 per cent of input cost. More productive labour alone raises the input bundle by 0.27 per cent. The rest depends on how much capital and intermediate inputs rise.
\item \textbf{The closure decides how much they rise.} With the foreign closure, capital and intermediate inputs rise 1.07 and 1.08 per cent. In the base model they rise 0.64 and 0.74 per cent, because the terms of trade fall and imported inputs do not become cheaper.
\item \textbf{Fixed costs decide the level.} Fixed costs are 11 per cent of the input bundle, and they do not change in the shock. Output therefore rises more than the inputs: 1.18 per cent for an input rise of 1.05 per cent.
\end{enumerate}

The test is a shock that also raises the fixed cost by one per cent.
With the foreign closure, GDP, private consumption, exports, and imports are then all 1.04 per cent above the baseline in 2075. In the base model, GDP is 0.88 per cent above.
The fixed cost therefore explains 0.25 of the 0.29 percentage points that GDP rises above one per cent with the closure.

\section{Scenarios}

All shocks start in 2026 and stay in all later years. Table~\ref{tab:scenarios} lists the four scenarios.
The productivity shock raises \code{aProd} by one per cent for all three labour types in all industries.
The fixed cost shock raises \code{qFixedCost\_i} by one per cent in all industries.
The baseline and the calibration of the closure are those of \code{foreign\_closure.tex}.

\begin{table}[H]
\centering
\small
\caption{Scenarios}
\label{tab:scenarios}
\begin{tabular}{@{}llll@{}}
\toprule
Short name & Model & Labour productivity & Fixed cost \\
\midrule
Closure, P & Base model with the foreign closure & $+1$ per cent & no change \\
Base, P & Base model & $+1$ per cent & no change \\
Closure, P+F & Base model with the foreign closure & $+1$ per cent & $+1$ per cent \\
Base, P+F & Base model & $+1$ per cent & $+1$ per cent \\
\bottomrule
\end{tabular}
\end{table}

All four models solve. The model with the closure has 2{,}885{,}340 equations, and the base model has 2{,}876{,}141.

\section{Output and inputs in the model}
\label{sec:block}

Each industry $i$ makes an input bundle $Q_{i,t}$ from capital, labour, and intermediate inputs. The bundle is a nested CES function. The elasticity is 0.7 in each of the nests \code{KE}, \code{KEL}, \code{KELB}, and \code{KELBM}.
Output $Y_{i,t}$ is what is left of the bundle after the fixed cost $F_{i,t}$:
\begin{equation}
Q_{i,t} - F_{i,t} = \theta_{i,t}\, Y_{i,t} .
\label{eq:top}
\end{equation}
In the code, $Q$ is \code{qProd} of the top nest, $F$ is \code{qFixedCost\_i}, and $\theta$ is \code{qTop2qY}. The fixed cost is exogenous and constant on the baseline.

Lower-case letters are per cent changes. To the first order, the change in the bundle is the sum of the input changes with cost shares as weights:
\begin{equation}
q = s_L\, (a + l) + s_K\, k + s_M\, m .
\label{eq:bundle}
\end{equation}
Here $a$ is labour productivity, $l$ is employment, $k$ is capital, and $m$ is intermediate inputs.
From (\ref{eq:top}), with $\varphi = F/Q$ the share of the fixed cost in the bundle and $f$ its per cent change,
\begin{equation}
y = \frac{q - \varphi f}{1 - \varphi} .
\label{eq:output}
\end{equation}
Table~\ref{tab:baseline} gives the shares on the baseline. With $\varphi = 0.110$, a rise of one per cent in the bundle with no change in the fixed cost raises output by 1.12 per cent.

\begin{table}[H]
\centering
\small
\caption{Shares on the baseline in 2075, all industries}
\label{tab:baseline}
\begin{tabular}{@{}lr@{}}
\toprule
 & Share \\
\midrule
Labour in input cost, $s_L$ & 0.274 \\
Capital in input cost, $s_K$ & 0.270 \\
Intermediate inputs in input cost, $s_M$ & 0.456 \\
Fixed cost in the input bundle, $\varphi$ & 0.110 \\
Imports in intermediate inputs & 0.378 \\
Imports in investment goods & 0.209 \\
\bottomrule
\end{tabular}
\end{table}

Real gross value added is output less intermediate inputs. It rises more than output when intermediate inputs rise less than output.

\section{Why GDP rises 1.29 per cent with the closure and 1.01 per cent in the base model}
\label{sec:why}

This section uses the two scenarios with the productivity shock only.
Table~\ref{tab:inputs} follows the shock from inputs to GDP, and Table~\ref{tab:contributions} applies (\ref{eq:bundle}).
The sum in Table~\ref{tab:contributions} equals the input bundle in Table~\ref{tab:inputs}, and (\ref{eq:output}) gives the row for output: $1.05 / 0.890 = 1.18$ and $0.78 / 0.890 = 0.88$.

\begin{table}[H]
\centering
\small
\caption{From inputs to GDP in 2075. Per cent deviation from the baseline}
\label{tab:inputs}
\begin{tabular}{@{}lrrrr@{}}
\toprule
 & Closure, P & Base, P & Closure, P+F & Base, P+F \\
\midrule
Employment & 0.00 & 0.00 & 0.00 & 0.00 \\
Labour productivity & 1.00 & 1.00 & 1.00 & 1.00 \\
Capital & 1.07 & 0.64 & 1.07 & 0.83 \\
Intermediate inputs & 1.08 & 0.74 & 1.03 & 0.78 \\
\midrule
Input bundle & 1.05 & 0.78 & 1.03 & 0.85 \\
Fixed cost & 0.00 & 0.00 & 1.00 & 1.00 \\
Output & 1.18 & 0.88 & 1.04 & 0.84 \\
\midrule
Real gross value added & 1.27 & 1.00 & 1.04 & 0.89 \\
Real GDP & 1.29 & 1.01 & 1.04 & 0.88 \\
\bottomrule
\end{tabular}
\end{table}

\begin{table}[H]
\centering
\small
\caption{Contributions to the input bundle in 2075, from (\ref{eq:bundle}). Percentage points}
\label{tab:contributions}
\begin{tabular}{@{}lrrrr@{}}
\toprule
 & Closure, P & Base, P & Closure, P+F & Base, P+F \\
\midrule
Labour & 0.27 & 0.27 & 0.27 & 0.27 \\
Capital & 0.29 & 0.17 & 0.29 & 0.22 \\
Intermediate inputs & 0.49 & 0.34 & 0.47 & 0.36 \\
\midrule
Sum & 1.05 & 0.78 & 1.03 & 0.85 \\
\bottomrule
\end{tabular}
\end{table}

\paragraph{Step 1: labour alone gives 0.27 per cent.}
Employment is the same in all four scenarios. The productivity shock raises the input bundle by $s_L \cdot 1 = 0.27$ per cent. All of the remaining rise comes from capital and intermediate inputs.

\paragraph{Step 2: the base model.}
Only the home economy becomes more productive. It must lower its price to sell more: the export price falls 0.16 per cent. Import prices are fixed, so the terms of trade fall 0.16 per cent (Table~\ref{tab:prices}).
Imports are 38 per cent of intermediate inputs and 21 per cent of investment goods. These inputs do not become cheaper when home prices fall.
The real wage rises 0.71 per cent. The wage rises 0.71 per cent relative to the user cost of capital and 0.68 per cent relative to the price of intermediate inputs.
Firms therefore add less capital and fewer intermediate inputs for each unit of labour: capital rises 0.64 per cent and intermediate inputs 0.74 per cent.

\paragraph{Step 3: the foreign closure.}
Foreign prices follow home prices, and foreign markets follow home imports. This is close to a rest of the world that gets the same rise in productivity.
The terms of trade do not change, and the prices of imported inputs follow the prices of home products.
The real wage rises 1.24 per cent, and the wage rises 1.47 per cent relative to the user cost of capital.
Capital and intermediate inputs rise 1.07 and 1.08 per cent, close to the rise of one per cent in labour in efficiency units.

\paragraph{Step 4: fixed costs.}
The fixed cost does not change, so (\ref{eq:output}) turns an input rise of 1.05 per cent into an output rise of 1.18 per cent. Intermediate inputs rise 1.08 per cent, less than output, so real gross value added rises 1.27 per cent.
The same step lifts the base model from 0.78 to 0.88 per cent in output and to 1.00 per cent in real gross value added. GDP in the base model is thus close to one per cent by coincidence: the fall in the terms of trade lowers it, and the fixed cost raises it.

\paragraph{Step 5: the ECB rate.}
With the closure, the ECB rate is 3.87 per cent in 2075, against 4.00 per cent. The hurdle rate for investment falls from 10.00 to 9.96 per cent, and the user cost of capital falls 0.31 per cent relative to the price of investment goods.
A rough calculation with the elasticity of 0.7 gives 0.16 per cent more capital and about 0.05 percentage points more GDP from this. It is not a model result.

\begin{table}[H]
\centering
\small
\caption{Prices and the ECB rate in 2075. Per cent deviation from the baseline. The two rates are in per cent}
\label{tab:prices}
\begin{tabular}{@{}lrrrr@{}}
\toprule
 & Closure, P & Base, P & Closure, P+F & Base, P+F \\
\midrule
Terms of trade & 0.00 & $-0.16$ & 0.00 & $-0.14$ \\
Consumer price & 0.00 & $-0.19$ & 0.00 & $-0.16$ \\
Real wage & 1.24 & 0.71 & 1.08 & 0.75 \\
Wage relative to the user cost of capital & 1.47 & 0.71 & 1.17 & 0.76 \\
Wage relative to the price of intermediate inputs & 1.19 & 0.68 & 1.06 & 0.73 \\
\midrule
ECB rate & 3.87 & 4.00 & 3.95 & 4.00 \\
Hurdle rate for investment & 9.96 & 10.00 & 9.98 & 10.00 \\
\bottomrule
\end{tabular}
\end{table}

\section{A shock to fixed costs also}

If the fixed cost is the cause of the rise above one per cent, a shock that also raises the fixed cost by one per cent must remove most of it. In (\ref{eq:output}), $f = 1$ and $q = 1$ give $y = 1$.

\subsection{Long run}

Table~\ref{tab:demand} gives the result in 2075.

\begin{table}[H]
\centering
\small
\caption{Quantities in 2075. Per cent deviation from the baseline}
\label{tab:demand}
\begin{tabular}{@{}lrrrr@{}}
\toprule
 & Closure, P & Base, P & Closure, P+F & Base, P+F \\
\midrule
GDP & 1.29 & 1.01 & 1.04 & 0.88 \\
Private consumption & 1.36 & 0.77 & 1.04 & 0.70 \\
Investment & 1.11 & 0.68 & 1.07 & 0.82 \\
Government consumption & 1.22 & 1.01 & 1.02 & 0.89 \\
Exports & 1.18 & 0.79 & 1.04 & 0.67 \\
Imports & 1.17 & 0.45 & 1.04 & 0.47 \\
\bottomrule
\end{tabular}
\end{table}

\begin{itemize}
\item \textbf{Foreign closure.} GDP is 1.04 per cent above the baseline, against 1.29 per cent with no shock to the fixed cost. GDP, private consumption, exports, and imports all rise 1.04 per cent. Investment rises 1.07 per cent and government consumption 1.02 per cent. The economy grows by close to one per cent in all parts.
\item \textbf{Base model.} GDP is 0.88 per cent above the baseline, against 1.01 per cent. The parts do not grow at the same rate: imports rise 0.47 per cent and exports 0.67 per cent. The terms of trade fall 0.14 per cent.
\item \textbf{Effect of the fixed cost.} It lowers GDP by 0.25 percentage points with the closure and by 0.12 points in the base model. With no change in the inputs, (\ref{eq:output}) gives $-0.12$ points on output. With the closure, the input bundle rises a little less (1.03 against 1.05 per cent). In the base model it rises more (0.85 against 0.78 per cent), and capital rises 0.83 per cent against 0.64.
\item \textbf{What is left.} With the closure, 0.04 percentage points above one per cent are left. The ECB rate is 3.95 per cent, against 4.00, and capital rises 1.07 per cent. Taxes on production and net pension saving are fixed amounts in the model and do not change in the shock. This note does not split the 0.04 points by cause.
\end{itemize}

\subsection{Paths}

Figure~\ref{fig:quantities} shows the quantities over time, Figure~\ref{fig:inputs} the inputs and output, and Figure~\ref{fig:prices} the prices and the ECB rate.
Blue lines are the foreign closure and orange lines are the base model. Solid lines have the productivity shock only. Dashed lines also have the fixed cost shock.

\begin{figure}[H]
\centering
\begin{tikzpicture}
\begin{groupplot}[fc, narrow, group style={group size=3 by 2, horizontal sep=1.0cm, vertical sep=1.4cm}]
\nextgroupplot[title={GDP}, ylabel={Per cent}, legend to name=fcflegend, legend columns=2]
\fourlines{gdp}
\fourlegend
\nextgroupplot[title={Private consumption}]
\fourlines{con}
\nextgroupplot[title={Investment}]
\fourlines{inv}
\nextgroupplot[title={Exports}, ylabel={Per cent}]
\fourlines{exp}
\nextgroupplot[title={Imports}]
\fourlines{imp}
\nextgroupplot[title={Government consumption}]
\fourlines{gov}
\end{groupplot}
\end{tikzpicture}

\ref{fcflegend}
\caption{Quantities. Per cent deviation from the baseline}
\label{fig:quantities}
\end{figure}

\begin{figure}[H]
\centering
\begin{tikzpicture}
\begin{groupplot}[fc, wide, group style={group size=2 by 2, horizontal sep=1.6cm, vertical sep=1.4cm}]
\nextgroupplot[title={Capital}, ylabel={Per cent}]
\fourlines{cap}
\nextgroupplot[title={Intermediate inputs}]
\fourlines{mat}
\nextgroupplot[title={Input bundle}, ylabel={Per cent}]
\fourlines{bun}
\nextgroupplot[title={Output}]
\fourlines{out}
\end{groupplot}
\end{tikzpicture}

\ref{fcflegend}
\caption{Inputs and output, sums over all industries at baseline prices. Per cent deviation from the baseline}
\label{fig:inputs}
\end{figure}

\begin{figure}[H]
\centering
\begin{tikzpicture}
\begin{groupplot}[fc, wide, group style={group size=2 by 2, horizontal sep=1.6cm, vertical sep=1.4cm}]
\nextgroupplot[title={Real wage}, ylabel={Per cent}]
\fourlines{rw}
\nextgroupplot[title={Terms of trade}]
\fourlines{tot}
\nextgroupplot[title={Consumer price}, ylabel={Per cent}]
\fourlines{pc}
\nextgroupplot[title={ECB rate, per cent}]
\fourlines{ecb}
\end{groupplot}
\end{tikzpicture}

\ref{fcflegend}
\caption{Prices and the ECB rate. Prices are in per cent deviation from the baseline}
\label{fig:prices}
\end{figure}

\FloatBarrier
\begin{itemize}
\item \textbf{First year.} Capital is slow, so all four scenarios start low. In 2026, GDP rises 0.59 and 0.62 per cent with the productivity shock, and 0.34 and 0.40 per cent with both shocks.
\item \textbf{Private consumption with the closure.} With both shocks it falls in the first years: $-0.27$ per cent in 2026 and $-0.37$ per cent in 2027. It is above the baseline from 2030. The ECB rate rises to 4.48 per cent in 2028 in this scenario, against 4.32 per cent with the productivity shock only.
\item \textbf{Investment.} In the base model, the fixed cost shock raises investment in the first years: 3.81 per cent in 2027, against 2.93 per cent. More fixed cost needs more inputs for the same output.
\item \textbf{End of the period.} With the closure and both shocks, all quantities in Table~\ref{tab:demand} change by less than 0.001 percentage points from 2064 to 2074.
\end{itemize}

\section{What the test shows}

\begin{itemize}
\item \textbf{With the closure, the model is close to homogeneous in productivity and fixed cost.} One per cent more of both gives about one per cent more of all quantities, with no change in the terms of trade. The rise of 1.29 per cent in GDP for the productivity shock alone comes from the fixed cost, which is 11 per cent of the input bundle and does not grow with the economy.
\item \textbf{In the base model, it is not.} One per cent more of both gives 0.88 per cent more GDP. The home economy alone becomes more productive, the terms of trade fall, and capital and intermediate inputs rise less than labour in efficiency units.
\item \textbf{The difference between the two models is the closure, not the fixed cost.} The gap in GDP is 0.29 percentage points with the productivity shock and 0.16 points with both shocks.
\end{itemize}

\section{Limits}

\begin{itemize}
\item \textbf{First-order accounting.} Equation (\ref{eq:bundle}) uses cost shares on the baseline. It fits the solved input bundle to two decimals in all four scenarios.
\item \textbf{Causes are read from the solution.} Steps 2 and 3 of Section~\ref{sec:why} compare prices and quantities in the solved scenarios. No scenario holds the terms of trade or the ECB rate fixed to isolate one cause. The effect of the ECB rate in step 5 is a rough calculation.
\item \textbf{Other items that do not grow.} The file \code{HomogeneityShock.jl} also scales net taxes on labour, other net production taxes, taxes and subsidies on production, net pension saving, and the foreign markets. This note did not. The 0.04 percentage points that are left with the closure can come from these items or from the ECB rate.
\item \textbf{Aggregates.} Capital, intermediate inputs, the input bundle, and output are sums over industries at baseline prices of the same year.
\item \textbf{Files.} The scenario files and the scripts that made the tables and the chart data are not in the repository. The chart data are in \code{docs/foreign\_closure\_fixed\_cost\_data.tex}.
\end{itemize}

\end{document}
